% !TEX root = ../main.tex
Il file \textit{stats\_utils.py} raccoglie le funzioni statistiche di base utilizzate nell'analisi dei dati sperimentali. Le funzioni supportano pesi opzionali $w_i = 1/\sigma_i^2$, utili quando le misure hanno incertezze diverse.

Le funzioni disponibili sono:
\begin{itemize}
    \item \texttt{median(x)}: calcola la mediana di un array di dati;
    \item \texttt{weighted\_mean(x, w)}: calcola la media aritmetica, opzionalmente pesata;
    \item \texttt{variance(x, w)}: calcola la varianza tramite la formula $\text{Var}(x) = E[x^2] - E[x]^2$;
    \item \texttt{covariance(x, y, w)}: calcola la covarianza tra due variabili tramite la formula $\text{Cov}(x,y) = E[xy] - E[x]E[y]$;
    \item \texttt{standard\_deviation(x, w)}: calcola la deviazione standard come $\sqrt{\text{Var}(x)}$.
\end{itemize}

\section{median(x)}
\label{sec:median}
\begin{apidesc}
  \item[\textbf{Scopo}]
    Calcolare la mediana di un array di dati numerici.
  \item[\textbf{Parametri}]\hfill
    \begin{description}[font=\normalfont\ttfamily\bfseries, labelwidth=1.5em, leftmargin=2em, itemsep=1pt, parsep=0pt, topsep=1pt]
      \item[x] array-like (lista, tupla, array NumPy) $\to$ convertito in \texttt{float64}
    \end{description}
  \item[\textbf{Vincoli}]
    Nessun vincolo aggiuntivo.
  \item[\textbf{Restituisce}]
    \texttt{float}: mediana di \texttt{x}.
  \item[\textbf{Eccezioni}]
    Nessuna eccezione esplicita; se \texttt{x} è vuoto, \texttt{np.median} restituisce
    \texttt{NaN} con un avviso di runtime.
\end{apidesc}

\paragraph{Implementazione}
\begin{minted}{python}
def median(x):
    x = np.asarray(x, dtype=float)
    return np.median(x)
\end{minted}
La funzione converte l'input in un array \texttt{float64} e delega il calcolo a \texttt{np.median}.

\paragraph{Significato matematico}
Data una sequenza di $n$ osservazioni ordinate $x_{(1)} \leq x_{(2)} \leq \cdots \leq x_{(n)}$, la mediana è il valore centrale:
$$
\text{med}(x) =
\begin{cases}
  x_{\left(\frac{n+1}{2}\right)} & \text{se } n \text{ è dispari} \\[2pt]
  \frac{x_{\left(\frac{n}{2}\right)} + x_{\left(\frac{n}{2}+1\right)}}{2} & \text{se } n \text{ è pari}
\end{cases}
$$
A differenza della media, la mediana è robusta rispetto ai valori anomali (outlier): un singolo dato molto distante dagli altri non la sposta significativamente.

\paragraph{Esempio d'uso}
\begin{minted}{python}
x = [3.1, 2.7, 4.0, 2.9, 3.5]
print(median(x))  # 3.1
# Ordinati: [2.7, 2.9, 3.1, 3.5, 4.0] -> elemento centrale: 3.1
\end{minted}

\newpage
\section{\texorpdfstring{weighted\_mean(x, w=None)}{weighted\_mean(x, w=None)}}
\label{sec:weighted-mean}
\begin{apidesc}
  \item[\textbf{Scopo}]
    Calcolare la media aritmetica di un array, con possibilità di pesare ogni dato
    in base alla sua incertezza.
  \item[\textbf{Parametri}]\hfill
    \begin{description}[font=\normalfont\ttfamily\bfseries, labelwidth=1.5em, leftmargin=2em, itemsep=1pt, parsep=0pt, topsep=1pt]
      \item[x] array-like $\to$ convertito in \texttt{float64}
      \item[w] array-like o \texttt{None} (default: \texttt{None})
    \end{description}
  \item[\textbf{Vincoli}]
    Se \texttt{w} $\neq$ \texttt{None}, i pesi devono essere pre-calcolati come
    $w_i = 1/\sigma_i^2$ e $\sum_i w_i > 0$. Se \texttt{w = None}, si usa
    $w_i = 1$ per tutti gli $i$.
  \item[\textbf{Restituisce}]
    \texttt{float}: media (pesata) di \texttt{x}.
  \item[\textbf{Eccezioni}]
    Nessuna eccezione esplicita; se $\sum_i w_i = 0$, NumPy esegue una divisione
    per zero floating-point e restituisce \texttt{NaN}.
\end{apidesc}

\paragraph{Implementazione}
\begin{minted}{python}
def weighted_mean(x, w=None):
    x = np.asarray(x, dtype=float)
    if w is None:
        w = np.ones(len(x))
    else:
        w = np.asarray(w, dtype=float)
    return np.sum(x * w) / np.sum(w)
\end{minted}
Se non si passa un valore \texttt{w}, la funzione crea un array di tutti \texttt{1},
quindi ogni dato pesa uguale e si ottiene la media aritmetica classica.

\paragraph{Significato matematico}
Senza pesi ($w_i = 1$ per tutti) si ha la media aritmetica classica:
$$\overline{x} = \frac{1}{N} \sum_i x_i$$

Se vengono passati i pesi $w_i = 1/\sigma_i^2$, si calcola la media pesata:
$$\overline{x}_w = \frac{\sum_i w_i x_i}{\sum_i w_i}$$
dove i punti con incertezza minore (peso maggiore) contribuiscono di più al risultato.

\paragraph{Esempio d'uso}
\begin{minted}{python}
import numpy as np
from stats_utils import weighted_mean

# Media semplice
x = [1.0, 2.0, 3.0, 4.0, 5.0]
print(weighted_mean(x))         # 3.0

# Media pesata: misure con incertezze diverse
x     = np.array([10.0, 10.5,  9.8])
sigma = np.array([ 0.2,  0.5,  0.1])
w     = 1 / sigma**2            # [25.0, 4.0, 100.0]
print(weighted_mean(x, w))      # ≈ 9.860 (dominato da x=9.8, sigma=0.1)
\end{minted}

\section{variance(x, w=None)}
\label{sec:variance}
\begin{apidesc}
  \item[\textbf{Scopo}]
    Calcolare la varianza di un array, con possibilità di pesare ogni dato.
  \item[\textbf{Parametri}]\hfill
    \begin{description}[font=\normalfont\ttfamily\bfseries, labelwidth=1.5em, leftmargin=2em, itemsep=1pt, parsep=0pt, topsep=1pt]
      \item[x] array-like $\to$ convertito in \texttt{float64}
      \item[w] array-like o \texttt{None} (default: \texttt{None})
    \end{description}
  \item[\textbf{Vincoli}]
    Se \texttt{w} $\neq$ \texttt{None}, i pesi devono essere pre-calcolati come
    $w_i = 1/\sigma_i^2$ e $\sum_i w_i > 0$. Se \texttt{w = None}, si usa
    $w_i = 1$ per tutti gli $i$.
  \item[\textbf{Restituisce}]
    \texttt{float}: varianza di \texttt{x}.
  \item[\textbf{Eccezioni}]
    Nessuna eccezione esplicita; le stesse condizioni limite di
    \texttt{weighted\_mean} si propagano (divisione per zero se $\sum_i w_i = 0$).
\end{apidesc}

\paragraph{Implementazione}
\begin{minted}{python}
def variance(x, w=None):
    x = np.asarray(x, dtype=float)
    return weighted_mean(x**2, w) - weighted_mean(x, w)**2
\end{minted}
La funzione delega interamente il calcolo a \texttt{weighted\_mean}, sfruttando
l'identità algebrica tra media dei quadrati e quadrato della media.

\paragraph{Significato matematico}
La varianza è calcolata tramite l'identità:
$$\text{Var}(x) = \overline{x^2} - \overline{x}^2$$
ovvero la media dei quadrati meno il quadrato della media. Con i pesi, la formula completa diventa:
$$\text{Var}_w(x) = \frac{\sum_i w_i x_i^2}{\sum_i w_i} - \left( \frac{\sum_i w_i x_i}{\sum_i w_i} \right)^2$$
Senza pesi ($w_i = 1$ per tutti) si riduce alla forma classica:
$$\text{Var}(x) = \frac{1}{N} \sum_i (x_i - \overline{x})^2$$

\begin{quote}
\textbf{Nota:} tutte le funzioni di questo modulo usano statistiche descrittive della
popolazione (divisore $N$), non stimatori campionari non distorti (divisore $N-1$).
La notazione $E[x^2] - E[x]^2$ si riferisce pertanto alla media empirica del campione,
non al valore atteso di una variabile aleatoria.
\end{quote}

\paragraph{Esempio d'uso}
\begin{minted}{python}
from stats_utils import variance

x = [1.0, 2.0, 3.0, 4.0, 5.0]
print(variance(x))   # 2.0
# media = 3, E[x^2] = (1+4+9+16+25)/5 = 11 -> 11 - 9 = 2
\end{minted}

\section{covariance(x, y, w=None)}
\label{sec:covariance}
\begin{apidesc}
  \item[\textbf{Scopo}]
    Calcolare la covarianza tra due array di dati, con possibilità di pesare ogni coppia.
  \item[\textbf{Parametri}]\hfill
    \begin{description}[font=\normalfont\ttfamily\bfseries, labelwidth=1.5em, leftmargin=2em, itemsep=1pt, parsep=0pt, topsep=1pt]
      \item[x] array-like $\to$ convertito in \texttt{float64}
      \item[y] array-like $\to$ convertito in \texttt{float64}
      \item[w] array-like o \texttt{None} (default: \texttt{None})
    \end{description}
  \item[\textbf{Vincoli}]
    \texttt{x} e \texttt{y} devono avere la stessa lunghezza.
    Se \texttt{w} $\neq$ \texttt{None}, i pesi devono essere pre-calcolati come
    $w_i = 1/\sigma_i^2$ e $\sum_i w_i > 0$.
  \item[\textbf{Restituisce}]
    \texttt{float}: covarianza tra \texttt{x} e \texttt{y}.
  \item[\textbf{Eccezioni}]
    \texttt{ValueError} se \texttt{len(x) != len(y)}.
\end{apidesc}

\paragraph{Implementazione}
\begin{minted}{python}
def covariance(x, y, w=None):
    x = np.asarray(x, dtype=float)
    y = np.asarray(y, dtype=float)
    if len(x) != len(y):
        raise ValueError("x and y must have the same length")
    return weighted_mean(x * y, w) - weighted_mean(x, w) * weighted_mean(y, w)
\end{minted}
Dopo aver verificato che \texttt{x} e \texttt{y} abbiano la stessa lunghezza,
la funzione sfrutta la stessa identità algebrica usata per la varianza, applicata
al prodotto $x_i y_i$.

\paragraph{Significato matematico}
La covarianza sfrutta l'identità:
$$\text{Cov}(x,y) = \overline{xy} - \overline{x}\,\overline{y}$$

I tre termini nel codice corrispondono a:
\begin{align*}
  \texttt{weighted\_mean(x * y, w)} &= \overline{xy}_w = \frac{\sum_i w_i \, x_i \, y_i}{\sum_i w_i} \\[4pt]
  \texttt{weighted\_mean(x, w) * weighted\_mean(y, w)} &= \bar{x}_w \, \bar{y}_w =
    \frac{\sum_i w_i \, x_i}{\sum_i w_i} \cdot \frac{\sum_i w_i \, y_i}{\sum_i w_i}
\end{align*}

Il risultato è equivalente alla forma pesata classica:
$$\text{Cov}_w(x,y) = \frac{\sum_i w_i (x_i - \bar{x}_w)(y_i - \bar{y}_w)}{\sum_i w_i}$$

Una covarianza positiva indica che $x$ e $y$ tendono a crescere insieme;
negativa, che crescono in direzioni opposte; zero, che sono linearmente indipendenti.

\paragraph{Esempio d'uso}
\begin{minted}{python}
from stats_utils import covariance

x = [1.0, 2.0, 3.0, 4.0]
y = [2.1, 3.9, 6.2, 7.8]   # y cresce con x -> covarianza positiva
print(covariance(x, y))     # ≈ 2.425
x2 = [1.0, 2.0, 3.0]
y2 = [1.0, 2.0, 3.0]
print(covariance(x2, y2))   # = variance(x2) ≈ 0.667
\end{minted}


\section{\texorpdfstring{standard\_deviation(x, w=None)}{standard\_deviation(x, w=None)}}
\label{sec:standard-deviation}
\begin{apidesc}
  \item[\textbf{Scopo}]
    Calcolare la deviazione standard di un array, con possibilità di pesare ogni dato.
  \item[\textbf{Parametri}]\hfill
    \begin{description}[font=\normalfont\ttfamily\bfseries, labelwidth=1.5em, leftmargin=2em, itemsep=1pt, parsep=0pt, topsep=1pt]
      \item[x] array-like $\to$ convertito in \texttt{float64}
      \item[w] array-like o \texttt{None} (default: \texttt{None})
    \end{description}
  \item[\textbf{Vincoli}]
    \texttt{x} non deve essere vuoto ($N \geq 1$).
    Se \texttt{w} $\neq$ \texttt{None}, i pesi devono essere pre-calcolati come
    $w_i = 1/\sigma_i^2$ e $\sum_i w_i > 0$.
  \item[\textbf{Restituisce}]
    \texttt{float}: deviazione standard di \texttt{x}.
  \item[\textbf{Eccezioni}]
    \texttt{ValueError} se \texttt{x} è vuoto ($N = 0$).
\end{apidesc}

\paragraph{Implementazione}
\begin{minted}{python}
def standard_deviation(x, w=None):
    x = np.asarray(x, dtype=float)
    if len(x) == 0:
        raise ValueError("x must contain at least one value")
    return np.sqrt(variance(x, w))
\end{minted}
Prima di procedere, la funzione controlla che \texttt{x} non sia vuoto,
poi restituisce la radice quadrata della varianza calcolata da \texttt{variance}.

\paragraph{Significato matematico}
La deviazione standard è la radice quadrata della varianza:
$$\sigma(x) = \sqrt{\text{Var}(x)}$$
Con i pesi la formula completa è:
$$\sigma_w(x) = \sqrt{\frac{\sum_i w_i x_i^2}{\sum_i w_i} - \left( \frac{\sum_i w_i x_i}{\sum_i w_i} \right)^2}$$
Senza pesi ($w_i = 1$ per tutti) si riduce alla forma classica:
$$\sigma(x) = \sqrt{\frac{1}{N} \sum_i (x_i - \overline{x})^2}$$
La deviazione standard ha le stesse unità di misura dei dati, a differenza della
varianza che è espressa in unità al quadrato: è quindi la grandezza tipicamente
riportata come incertezza statistica di una serie di misure.

\paragraph{Esempio d'uso}
\begin{minted}{python}
from stats_utils import standard_deviation

x = [2.0, 4.0, 4.0, 4.0, 5.0, 5.0, 7.0, 9.0]
print(standard_deviation(x))   # 2.0
# media = 5, varianza = 4 -> sqrt(4) = 2
\end{minted}
